\documentclass{amsart}
\usepackage{amsmath,amssymb,amsthm,hyperref}
\newtheorem{theorem}{Theorem}
\newtheorem{lemma}{Lemma}
\newtheorem{proposition}{Proposition}
\newtheorem{corollary}{Corollary}
\theoremstyle{definition}
\newtheorem{definition}{Definition}
\newtheorem{assumption}{Assumption}
\newtheorem{remark}{Remark}

\begin{document}

\title{A Mediation Model Linking Team Efficacy, Learning Behavior, and Team Performance}
\maketitle

\section{Formalization}

The problem is a substantive (organizational--behavioral) claim. To
\emph{justify} it rigorously we recast the three latent constructs as
real-valued random variables defined on a common probability space and we make
the causal claim precise as a statement about a linear structural causal model
(SCM). ``Partial mediation'' is then a precise, checkable algebraic condition,
and ``conditions under which the mediation holds'' becomes a precise list of
identifying assumptions. The model used is the standard linear mediation /
path-analytic model formalized by Wright's path analysis and made operational
by the Baron--Kenny three-equation framework and its causal-inference
refinement; we prove every identity and every inequality from the axioms below
rather than citing them.

\begin{definition}[Population of teams]
Fix a population of work teams and a probability space
$(\Omega,\mathcal F,\mathbb P)$ whose sample points $\omega\in\Omega$ index
teams drawn at random from this population. All random variables below are
square-integrable functions $\Omega\to\mathbb R$, i.e.\ elements of
$L^2(\Omega,\mathcal F,\mathbb P)$.
\end{definition}

\begin{definition}[Constructs as random variables]\label{def:constructs}
For a randomly drawn team $\omega$ define:
\begin{itemize}
\item $E(\omega)$ \emph{(team efficacy)}: a scalar score for the shared belief
in the team's collective capability, e.g.\ the mean over members of a validated
collective-efficacy item battery.
\item $L(\omega)$ \emph{(learning behavior)}: a scalar score for the team's
information-acquisition/processing actions (feedback seeking, help seeking,
error discussion, experimentation), e.g.\ a behavioral or survey index.
\item $P(\omega)$ \emph{(team performance)}: a scalar for goal accomplishment,
e.g.\ a supervisor rating or objective output index.
\end{itemize}
Without loss of generality we center all variables:
$\mathbb E[E]=\mathbb E[L]=\mathbb E[P]=0$ (replace each variable by its
deviation from its mean; this changes no covariance or regression slope). We
assume $\operatorname{Var}(E)>0$ and $\operatorname{Var}(L)>0$, i.e.\ teams
genuinely differ in efficacy and in learning behavior.
\end{definition}

\begin{definition}[Linear structural causal model]\label{def:scm}
The substantive theory ``efficacy raises confidence to learn, and both efficacy
and learning drive performance'' is encoded by the structural equations
\begin{align}
L &= a\,E + U_L, \label{eq:struct-L}\\
P &= c'\,E + b\,L + U_P, \label{eq:struct-P}
\end{align}
with structural coefficients $a,b,c'\in\mathbb R$ and disturbances
$U_L,U_P\in L^2$. Equations \eqref{eq:struct-L}--\eqref{eq:struct-P} are
\emph{structural}: each is read as ``the left side is generated by the right
side,'' so that an external (experimental) intervention $\mathrm{do}(E=e)$
replaces $E$ by the constant $e$ on the right-hand sides while leaving
$a,b,c',U_L,U_P$ unchanged (Pearl's $\mathrm{do}$-calculus reading of an SCM).
\end{definition}

\begin{assumption}[Exogeneity / no unmeasured confounding]\label{ass:exo}
The disturbances are mean-zero and orthogonal to the regressors that precede
them in the causal order:
\[
\mathbb E[U_L]=\mathbb E[U_P]=0,\qquad
\mathbb E[E\,U_L]=0,\qquad
\mathbb E[E\,U_P]=0,\qquad
\mathbb E[L\,U_P]=0 .
\]
Substantively: there is no third variable that jointly causes $E$ and $L$
(beyond what is modeled), none that jointly causes $E$ and $P$, and none that
jointly causes $L$ and $P$. (This is the identifying ``no-confounding''
condition; under randomized assignment of $E$ it holds by design for the
$E$-orthogonality conditions, and sequential ignorability gives the
$L$-orthogonality condition.)
\end{assumption}

\begin{remark}[What ``partial mediation'' means]
Write the \emph{total effect} of $E$ on $P$ as the coefficient $c$ in the
population linear projection of $P$ on $E$, the \emph{direct effect} as $c'$,
and the \emph{indirect (mediated) effect} as the product $ab$. The theory's
claim ``learning behavior \emph{partially} mediates the efficacy--performance
relation'' is the conjunction
\[
\underbrace{ab\neq 0}_{\text{a real indirect path}}
\quad\text{and}\quad
\underbrace{c'\neq 0}_{\text{efficacy still acts directly}} ,
\]
together with the bookkeeping identity $c=c'+ab$. \emph{Full} mediation is the
boundary case $c'=0$; \emph{no} mediation is $ab=0$. Below we prove the identity
unconditionally and then give exact sign/magnitude conditions for partial
mediation.
\end{remark}

\section{Main results}

\subsection{Identification of the path coefficients}

\begin{lemma}[Identification]\label{lem:ident}
Under Definition~\ref{def:scm} and Assumption~\ref{ass:exo}, the structural
coefficients equal population regression slopes:
\[
a=\frac{\operatorname{Cov}(E,L)}{\operatorname{Var}(E)},\qquad
\begin{pmatrix} c'\\ b\end{pmatrix}
=\Sigma^{-1}
\begin{pmatrix}\operatorname{Cov}(E,P)\\ \operatorname{Cov}(L,P)\end{pmatrix},
\quad
\Sigma=\begin{pmatrix}\operatorname{Var}(E)&\operatorname{Cov}(E,L)\\
\operatorname{Cov}(E,L)&\operatorname{Var}(L)\end{pmatrix},
\]
provided $\Sigma$ is invertible, i.e.\ $E$ and $L$ are not perfectly
collinear.
\end{lemma}

\begin{proof}
\emph{Coefficient $a$.} Multiply \eqref{eq:struct-L} by $E$ and take
expectations:
\[
\mathbb E[EL]=a\,\mathbb E[E^2]+\mathbb E[E\,U_L].
\]
By centering (Definition~\ref{def:constructs}), $\mathbb E[EL]=\operatorname{Cov}(E,L)$
and $\mathbb E[E^2]=\operatorname{Var}(E)$; by Assumption~\ref{ass:exo},
$\mathbb E[E\,U_L]=0$. Hence $\operatorname{Cov}(E,L)=a\operatorname{Var}(E)$, and
since $\operatorname{Var}(E)>0$ we may divide to get the stated formula.

\emph{Coefficients $(c',b)$.} Multiply \eqref{eq:struct-P} successively by $E$
and by $L$ and take expectations:
\begin{align}
\mathbb E[EP]&=c'\,\mathbb E[E^2]+b\,\mathbb E[EL]+\mathbb E[E\,U_P],\label{eq:nE}\\
\mathbb E[LP]&=c'\,\mathbb E[EL]+b\,\mathbb E[L^2]+\mathbb E[L\,U_P].\label{eq:nL}
\end{align}
By centering, the raw moments are the covariances/variances in $\Sigma$ and in
the right-hand vector. By Assumption~\ref{ass:exo}, $\mathbb E[E\,U_P]=0$ and
$\mathbb E[L\,U_P]=0$. Thus \eqref{eq:nE}--\eqref{eq:nL} read
\[
\begin{pmatrix}\operatorname{Cov}(E,P)\\ \operatorname{Cov}(L,P)\end{pmatrix}
=\Sigma\begin{pmatrix}c'\\ b\end{pmatrix}.
\]
Because $\Sigma$ is a $2\times2$ covariance matrix it is positive
semidefinite; the hypothesis that $E,L$ are not perfectly collinear means the
inequality in Cauchy--Schwarz,
$\operatorname{Cov}(E,L)^2\le\operatorname{Var}(E)\operatorname{Var}(L)$, is
strict, so $\det\Sigma>0$ and $\Sigma$ is invertible. Left-multiplying by
$\Sigma^{-1}$ gives the formula.
\end{proof}

\subsection{The exact effect decomposition}

\begin{theorem}[Total $=$ direct $+$ indirect]\label{thm:decomp}
Define the total effect $c:=\operatorname{Cov}(E,P)/\operatorname{Var}(E)$, the
population slope of $P$ regressed on $E$ alone. Under
Definition~\ref{def:scm} and Assumption~\ref{ass:exo},
\begin{equation}\label{eq:decomp}
\boxed{\,c=c'+ab\,}.
\end{equation}
Moreover $c$ also equals the average causal effect:
$c=\dfrac{d}{de}\,\mathbb E[P\mid \mathrm{do}(E=e)]$, and $ab$ equals the part
of that effect transmitted through $L$.
\end{theorem}

\begin{proof}
\emph{Algebraic identity.} Substitute the structural equation
\eqref{eq:struct-L} for $L$ into \eqref{eq:struct-P}:
\[
P=c'E+b(aE+U_L)+U_P=(c'+ab)E+(bU_L+U_P).
\]
Take covariance with $E$:
\[
\operatorname{Cov}(E,P)=(c'+ab)\operatorname{Var}(E)
+b\,\mathbb E[E\,U_L]+\mathbb E[E\,U_P].
\]
By Assumption~\ref{ass:exo} the last two terms vanish, so
$\operatorname{Cov}(E,P)=(c'+ab)\operatorname{Var}(E)$. Dividing by
$\operatorname{Var}(E)>0$ and using the definition of $c$ yields
$c=c'+ab$.

\emph{Causal interpretation.} Apply $\mathrm{do}(E=e)$: by
Definition~\ref{def:scm} this sets $E=e$ on the right-hand sides while keeping
$a,b,c',U_L,U_P$ fixed, giving the post-intervention variables
$L^{(e)}=ae+U_L$ and
$P^{(e)}=c'e+bL^{(e)}+U_P=(c'+ab)e+(bU_L+U_P)$. Taking expectations and using
$\mathbb E[U_L]=\mathbb E[U_P]=0$,
\[
\mathbb E[P\mid\mathrm{do}(E=e)]=\mathbb E[P^{(e)}]=(c'+ab)\,e,
\]
whose derivative in $e$ is $c'+ab=c$. The portion routed through $L$ is the
$\mathrm{do}$-effect of $E$ on $L$, namely $\frac{d}{de}\mathbb E[L\mid\mathrm{do}(E=e)]=a$,
composed with the $\mathrm{do}$-effect of $L$ on $P$ holding $E$ fixed, which by
\eqref{eq:struct-P} is $b$; their product $ab$ is the indirect effect.
\end{proof}

\begin{remark}[Faithfulness check]
Identity \eqref{eq:decomp} is exactly the bookkeeping behind the Baron--Kenny
procedure: regress $P$ on $E$ to get $c$; regress $L$ on $E$ to get $a$;
regress $P$ on $E$ and $L$ to get $c'$ and $b$. We verified \eqref{eq:decomp}
numerically (non-Gaussian disturbances) and recovered $c=c'+ab$ to floating
precision, confirming Gaussianity is \emph{not} required---only the second-order
orthogonality of Assumption~\ref{ass:exo}.
\end{remark}

\subsection{Conditions for partial mediation}

\begin{theorem}[Partial-mediation criterion]\label{thm:partial}
Adopt Definition~\ref{def:scm} and Assumption~\ref{ass:exo}, with $E,L$ not
perfectly collinear. Then learning behavior \emph{partially mediates} the
efficacy--performance relationship---defined (as in the Remark above) by the
conjunction of a nonzero mediated component $ab\neq0$ and a nonzero residual
direct component $c'\neq0$ in the decomposition $c=c'+ab$---\emph{if and only
if}
\begin{equation}\label{eq:pm-cond}
a\neq0,\qquad b\neq0,\qquad c'\neq0 .
\end{equation}
Under \eqref{eq:pm-cond} the indirect effect $ab$ and the direct effect $c'$
are both nonzero, and the total effect $c=c'+ab$ is nonzero whenever
$c'$ and $ab$ do not exactly cancel, i.e.\ whenever $c'\neq-ab$.
\end{theorem}

\begin{proof}
By Theorem~\ref{thm:decomp} the decomposition $c=c'+ab$ holds, so the indirect
component is the single number $ab$ and the direct component is $c'$.

($\Leftarrow$) If $a\neq0$ and $b\neq0$ then $ab\neq0$, so the mediated path
carries a genuine effect; if additionally $c'\neq0$ the direct path is also
nonzero. Both components being nonzero is exactly partial mediation.

($\Rightarrow$) Conversely suppose mediation is partial, i.e.\ $ab\neq0$ and
$c'\neq0$. From $ab\neq0$ we get $a\neq0$ and $b\neq0$ (a product is nonzero iff
both factors are). Together with $c'\neq0$ this is \eqref{eq:pm-cond}.

Finally, $c=c'+ab$ is zero only if $c'=-ab$; outside this single cancellation
hyperplane, $c\neq0$, so a marginal efficacy--performance association is
observable.
\end{proof}

\begin{corollary}[Sign and dominance of the mediated path]\label{cor:sign}
Under \eqref{eq:pm-cond}:
\begin{enumerate}
\item \emph{(Positive theory.)} If efficacy raises learning ($a>0$) and learning
raises performance ($b>0$), then the indirect effect $ab>0$: shared efficacy
beliefs improve performance \emph{in part through} increased learning behavior.
\item \emph{(Proportion mediated.)} Suppose the two routes act in the same
direction, i.e.\ $\operatorname{sign}(ab)=\operatorname{sign}(c')$ (both
positive or both negative). Then $c=c'+ab$ has that same sign and is nonzero,
and the fraction of the total effect transmitted through learning behavior,
\[
\Pi:=\frac{ab}{c}=\frac{ab}{c'+ab},
\]
satisfies $\Pi\in(0,1)$, strictly between $0$ (no mediation) and $1$ (full
mediation); moreover $\Pi$ is strictly increasing in $|ab|$ and strictly
decreasing in $|c'|$.
\item \emph{(Full vs.\ no mediation as limits.)} With $ab$ and $c'$ of the same
sign held, $c'\to0$ gives $\Pi\to1$ (full mediation) and $ab\to0$ gives
$\Pi\to0$ (no mediation), so partial mediation is precisely the interior
regime $\Pi\in(0,1)$.
\end{enumerate}
\end{corollary}

\begin{proof}
(1) The product of two positive reals is positive, so $a,b>0\Rightarrow ab>0$;
by Theorem~\ref{thm:decomp} this $ab$ is the performance gain routed through
$L$.

(2) Write $s=\operatorname{sign}(ab)=\operatorname{sign}(c')\in\{+1,-1\}$, so
$ab=s|ab|$ and $c'=s|c'|$ with $|ab|>0$, $|c'|>0$ by \eqref{eq:pm-cond}. Then
\[
c=c'+ab=s\bigl(|c'|+|ab|\bigr),
\]
which has sign $s$ and magnitude $|c'|+|ab|>0$, so $c\neq0$. Hence
\[
\Pi=\frac{ab}{c}=\frac{s|ab|}{s(|c'|+|ab|)}=\frac{|ab|}{|c'|+|ab|}.
\]
Since $|ab|>0$ and $|c'|>0$, the numerator is positive and strictly smaller
than the denominator, so $0<\Pi<1$. For monotonicity, set $x=|ab|>0$,
$y=|c'|>0$ and view $\Pi=x/(x+y)$; then
\[
\frac{\partial\Pi}{\partial x}=\frac{y}{(x+y)^2}>0,\qquad
\frac{\partial\Pi}{\partial y}=\frac{-x}{(x+y)^2}<0,
\]
so $\Pi$ strictly increases in $|ab|$ and strictly decreases in $|c'|$.

(3) From $\Pi=|ab|/(|c'|+|ab|)$: as $|c'|\to0^+$, $\Pi\to|ab|/|ab|=1$; as
$|ab|\to0^+$, $\Pi\to0/|c'|=0$. Thus the open interval $\Pi\in(0,1)$ is exactly
the partial-mediation regime, with the two endpoints corresponding to full and
no mediation.
\end{proof}

\begin{remark}[Why same-sign routes are needed for $\Pi\in(0,1)$]
If $ab$ and $c'$ had opposite signs (a \emph{suppression} or inconsistent
mediation pattern), $\Pi=ab/c$ can exceed $1$ or be negative; e.g.\ $ab=2$,
$c'=-1/2$ give $c=3/2$ and $\Pi=4/3>1$. The hypothesis
$\operatorname{sign}(ab)=\operatorname{sign}(c')$ in
Corollary~\ref{cor:sign}(2) rules out exactly this case and is what guarantees
the natural ``proportion mediated'' interpretation $\Pi\in(0,1)$.
\end{remark}

\subsection{Mechanism justification (the substantive claim)}

\begin{proposition}[Why efficacy aids performance, and the role of learning]
\label{prop:mech}
Under Definition~\ref{def:scm} and Assumption~\ref{ass:exo}, the total benefit
of higher team efficacy for performance, $c$, operates through exactly two
mechanisms, and only these two, and they sum to $c$:
\begin{enumerate}
\item a \emph{direct} mechanism of magnitude $c'$ (efficacy raises performance
through effort, persistence, goal-setting, and motivational pathways not routed
through learning behavior), and
\item an \emph{indirect, learning-mediated} mechanism of magnitude $ab$, itself
the composition of (i) efficacy raising members' confidence to engage in
learning behavior, magnitude $a$, and (ii) learning behavior raising
performance, magnitude $b$.
\end{enumerate}
The substantive claim ``team efficacy improves performance, at least in part,
by raising confidence to engage in learning behavior'' is true precisely under
condition \eqref{eq:pm-cond} together with $a>0,b>0$ (so that the learning path
is real and performance-enhancing), and it is then quantified by $ab>0$ and,
when additionally $c'>0$, by the mediated proportion $\Pi=ab/c\in(0,1)$.
\end{proposition}

\begin{proof}
Substituting \eqref{eq:struct-L} into \eqref{eq:struct-P} (as in the proof of
Theorem~\ref{thm:decomp}) gives the reduced form
$P=(c'+ab)E+(bU_L+U_P)$. The coefficient on $E$ is $c'+ab$, an additive sum of
exactly the two terms $c'$ and $ab$; the SCM in Definition~\ref{def:scm} admits
no other directed path from $E$ to $P$ (the only paths in the path diagram
$E\!\to\!P$ and $E\!\to\!L\!\to\!P$ are the direct edge and the two-edge chain
through $L$), so by the path-tracing rules for linear SCMs these are the only
two mechanisms, establishing exhaustiveness. Their magnitudes $c'$ and $ab$
were identified in Lemma~\ref{lem:ident} and Theorem~\ref{thm:decomp}. The
truth conditions for the verbal claim follow from Theorem~\ref{thm:partial} and
Corollary~\ref{cor:sign}: $a>0$ encodes ``efficacy raises confidence to learn,''
$b>0$ encodes ``learning raises performance,'' their product $ab>0$ is the
mediated benefit, and $c'\neq0$ makes the mediation partial rather than full;
when further $c'>0$ both routes reinforce and $\Pi=ab/c\in(0,1)$.
\end{proof}

\section{Summary of the model and its conditions}

\begin{itemize}
\item \textbf{Model.} Two structural equations \eqref{eq:struct-L},
\eqref{eq:struct-P} with path diagram $E\to L\to P$ and $E\to P$; coefficients
$a$ (efficacy$\to$learning), $b$ (learning$\to$performance), $c'$ (direct
efficacy$\to$performance).
\item \textbf{Core identity.} $c=c'+ab$ (Theorem~\ref{thm:decomp}): the total
efficacy effect equals direct plus indirect.
\item \textbf{Conditions for partial mediation.} $a\neq0$, $b\neq0$, $c'\neq0$
(Theorem~\ref{thm:partial}); for a \emph{positive} learning mechanism add
$a>0,b>0$, and for the proportion-mediated interpretation $\Pi\in(0,1)$ add
$c'>0$ (Corollary~\ref{cor:sign}, Proposition~\ref{prop:mech}).
\item \textbf{Identifying assumptions.} Linearity (Def.~\ref{def:scm}),
positive variances and no perfect $E$--$L$ collinearity
(Def.~\ref{def:constructs}, Lem.~\ref{lem:ident}), and the no-confounding /
exogeneity orthogonality conditions (Assumption~\ref{ass:exo}; guaranteed by
randomization of $E$ plus sequential ignorability of $L$).
\item \textbf{Conclusion.} Under these conditions, learning behavior partially
mediates the efficacy--performance relationship, with mediated effect $ab>0$ and
(when routes share sign) mediated proportion $\Pi=ab/c\in(0,1)$.
\end{itemize}

\end{document}
